
\begin{proposition}[All rank-four monodromy stratum signatures]\label{cdp:st:C.2.UN.11}
Let $f:U\to\Delta$ be a proper holomorphic map from a
smooth complex threefold, with connected fibres and
smooth fibres complex two-tori away from zero. Take
an embedded resolution, unchanged away from the
central fibre, and write its full central divisor as
$\sum_jN_jE_j$, with simple normal crossings. Put
\[
 E_j^\circ=E_j\setminus\bigcup_{i\ne j}E_i,
 \qquad a_n=\sum_{N_j=n}e(E_j^\circ).
 \tag{C.MS.1}\label{cdp:eq:C.MS.1}
\]
Let $T$ be the actual monodromy on $H^1(F,\Z)$,
and let $M$ be the order of its semisimple part.
Then $T\in\operatorname{SL}_4(\Z)$ is
quasi-unipotent, and the following exact maps
determine the entire sequence $(a_n)$:
\[
 \begin{aligned}
 L_k&=\det(I-T^k)
       =\sum_{n\mid k}n a_n,\\
 a_n&=\frac1n\sum_{d\mid n}\mu(n/d)\det(I-T^d).
 \end{aligned}                                  \tag{C.MS.2}\label{cdp:eq:C.MS.2}
\]
Here $\mu$ is the integer M\"obius function.
The complete table follows; $\Phi_d$ denotes
the $d$th cyclotomic polynomial and every
unlisted $a_n$ is zero.

\begin{center}
\small
\begin{tabular}{ccl}
$P_T$ & $M$ & Nonzero $a_n$\\\hline
$\Phi_{1}^{4}$ & 1 & $\text{all zero}$ \\
$\Phi_{1}^{2}\Phi_{2}^{2}$ & 2 & $\text{all zero}$ \\
$\Phi_{1}^{2}\Phi_{3}$ & 3 & $\text{all zero}$ \\
$\Phi_{1}^{2}\Phi_{4}$ & 4 & $\text{all zero}$ \\
$\Phi_{1}^{2}\Phi_{6}$ & 6 & $\text{all zero}$ \\
$\Phi_{2}^{4}$ & 2 & $a_{1}=16,\quad a_{2}=-8$ \\
$\Phi_{3}^{2}$ & 3 & $a_{1}=9,\quad a_{3}=-3$ \\
$\Phi_{2}^{2}\Phi_{4}$ & 4 & $a_{1}=8,\quad a_{2}=-4$ \\
$\Phi_{4}^{2}$ & 4 & $a_{1}=4,\quad a_{2}=6,\quad a_{4}=-4$ \\
$\Phi_{5}$ & 5 & $a_{1}=5,\quad a_{5}=-1$ \\
$\Phi_{2}^{2}\Phi_{3}$ & 6 & $a_{1}=12,\quad a_{2}=-6,\quad a_{3}=-4,\quad a_{6}=2$ \\
$\Phi_{2}^{2}\Phi_{6}$ & 6 & $a_{1}=4,\quad a_{2}=-2,\quad a_{3}=4,\quad a_{6}=-2$ \\
$\Phi_{3}\Phi_{6}$ & 6 & $a_{1}=3,\quad a_{2}=3,\quad a_{3}=-1,\quad a_{6}=-1$ \\
$\Phi_{6}^{2}$ & 6 & $a_{1}=1,\quad a_{2}=4,\quad a_{3}=5,\quad a_{6}=-4$ \\
$\Phi_{8}$ & 8 & $a_{1}=2,\quad a_{2}=1,\quad a_{4}=3,\quad a_{8}=-2$ \\
$\Phi_{10}$ & 10 & $a_{1}=1,\quad a_{2}=2,\quad a_{5}=3,\quad a_{10}=-2$ \\
$\Phi_{3}\Phi_{4}$ & 12 & $a_{1}=6,\quad a_{2}=3,\quad a_{3}=-2,\quad a_{4}=-3,\quad a_{6}=-1,\quad a_{12}=1$ \\
$\Phi_{4}\Phi_{6}$ & 12 & $a_{1}=2,\quad a_{2}=5,\quad a_{3}=2,\quad a_{4}=-3,\quad a_{6}=-3,\quad a_{12}=1$ \\
$\Phi_{12}$ & 12 & $a_{1}=1,\quad a_{3}=1,\quad a_{4}=2,\quad a_{6}=2,\quad a_{12}=-2$ \\
\end{tabular}
\end{center}


In particular $a_n=0$ unless $n\mid M$, where
\[
                 M\in\{1,2,3,4,5,6,8,10,12\}.
 \tag{C.MS.3}\label{cdp:eq:C.MS.3}
\]
On $H^1(F,\Q)$ one also has
\[
                  (T^M-I)^2=0,
                  \qquad\operatorname{rank}(T^M-I)\le2.
 \tag{C.MS.10}\label{cdp:eq:C.MS.10}
\]
All $a_n$ vanish exactly when $1$ is an eigenvalue
of $T$. The table is the complete range of \eqref{cdp:eq:C.MS.2}
on orientation-preserving quasi-unipotent integral
rank-four matrices. Every actual filling has one
of these signatures. Every row is realized by a
map from a smooth projective threefold to $\PP^1$
whose general fibre is an abelian surface and
whose every fibre has Euler number zero. It
can be chosen to have exactly two singular
fibres, both with reducible simple-normal-crossings
support.

There are two useful consequences for the original
fibre before resolution. Write it as $\sum_i m_iS_i$,
put $m_0=\min_i m_i$, and let $S_i^\circ$ be the
smooth locus of $S_i$ with all its intersections
with the other components removed. Then
\[
 m_0\sum_{m_i=m_0}e(S_i^\circ)
                    =\det(I-T^{m_0})\ge0.        \tag{C.MS.4}\label{cdp:eq:C.MS.4}
\]
If every component has multiplicity greater than
one, the entire signature is zero. For an
irreducible full fibre $mS$ one has
\[
 \begin{gathered}
 m>1\quad\Longrightarrow\quad e(S_{\rm reg})=0,\\
 e(S_{\rm reg})\in
       \{0,1,2,3,4,5,6,8,9,12,16\},\\
 e(S_{\rm reg})>0\quad\Longrightarrow\quad m=1.
 \end{gathered}                                  \tag{C.MS.5}\label{cdp:eq:C.MS.5}
\]
These constraints apply in particular to every
reducible or irreducible zero-Euler filling.
\end{proposition}
\begin{proof}
We specify first the comparison with the analytic
family. The resolution is an isomorphism over
$\Delta^*$, so it identifies the actual nearby
torus and its monodromy before and after resolution.
The proper global construction in
\cite[\S2, pp.~237--243, and \S3]{ACampo1975}
gives the nearby-cycle cohomology sheaves on the
resolved central support and their spectral sequence
to $H^*(F,\Q)$. In a chart
$z_1^{N_1}\cdots z_r^{N_r}=t$, a local nearby
fibre has $\gcd(N_1,\ldots,N_r)$ components,
each with the homotopy type of $(S^1)^{r-1}$.
Monodromy permutes these components cyclically;
after their permutation is killed its action
is a translation on the torus and induces the
identity on cohomology. A common multiple of all
the $N_j$ therefore acts identically on every
stalk of every nearby-cycle cohomology sheaf.
It acts identically on the corresponding
cohomology groups of those sheaves and hence
on the associated graded of the finite
spectral-sequence filtration of $H^*(F,\Q)$.
Its difference from the identity is nilpotent
on that cohomology. This proves quasi-unipotence
using the actual analytic normal-crossings
comparison; no ambient K\"ahler assumption enters.
In degree one this filtration has only the
two possible terms $H^1(E,\mathcal H^0\psi)$
and $H^0(E,\mathcal H^1\psi)$, where
$E=\bigcup_jE_j$ and $\mathcal H^q\psi$
denotes the nearby-cycle cohomology sheaf
already described. Thus for a
common multiplicity multiple $N$ one has
$(T^N-I)^2=0$. The semisimple order $M$
divides $N$, and $U=T^M$ is unipotent.
Writing
$U^{N/M}-I=(U-I)(I+U+\cdots+U^{N/M-1})$,
the second factor is invertible over $\Q$
and commutes with $U$. Hence $(U-I)^2=0$.
Its image is contained in its kernel on
the four-dimensional space, proving
the rank bound in \eqref{cdp:eq:C.MS.10}.

The same local model computes the alternating
stalk trace of the $k$th monodromy power. At a
point of $E_j^\circ$ it is $N_j$ if $N_j\mid k$
and zero otherwise. At a point lying on two or
more components it is zero, since the torus
factor has positive dimension and Euler number
zero. Taking the alternating trace through the
constructible spectral sequence, or its finite
cell filtration, gives
\[
 \sum_q(-1)^q\operatorname{Tr}
             (T^k;H^q(F,\Q))
       =\sum_{N_j\mid k}N_j e(E_j^\circ).
 \tag{C.MS.6}\label{cdp:eq:C.MS.6}
\]
This is the global form of the explicit local
calculation in \cite[Theorem~1, p.~235,
and proof on pp.~243--244]{ACampo1975}.
The cup-product map
$\Lambda^qH^1(F,\Z)\to H^q(F,\Z)$ is an
isomorphism for a torus and commutes with the
actual monodromy. Thus the left side of
\eqref{cdp:eq:C.MS.6} is exactly $\det(I-T^k)$.
The complex orientation of the smooth fibre
is preserved by transport. The map on its
top cohomology is consequently the identity;
under the fourth exterior power this says
$\det T=1$. Integer M\"obius inversion now
gives both equations of \eqref{cdp:eq:C.MS.2}.

We compute their complete range. The characteristic
polynomial of a quasi-unipotent integral rank-four
matrix is a product of cyclotomic polynomials
whose degrees sum to four. The possible individual
orders are
$1,2,3,4,5,6,8,10,12$. Since $\det T=1$,
the multiplicity of the eigenvalue $-1$ is even.
There are exactly nineteen resulting products:
the five containing $\Phi_1$ and the fourteen
remaining products displayed in the table.
This enumeration follows directly from
$\deg\Phi_1=\deg\Phi_2=1$,
$\deg\Phi_3=\deg\Phi_4=\deg\Phi_6=2$,
and degree four for $\Phi_5,\Phi_8,\Phi_{10},
\Phi_{12}$. The sum of the degree-one
multiplicities is even, so the $\Phi_1$
multiplicity is even as well.

The determinants $\det(I-T^k)$ depend only on
these eigenvalues, including their multiplicities;
Jordan blocks make no change to them. Each listed
product is represented by the block diagonal
matrix of its integral cyclotomic companion
matrices, which has determinant one. Its traces
therefore compute the exact full range on
matrices. Applying \eqref{cdp:eq:C.MS.2} gives the displayed
integers.

There is also a direct verification that no
unlisted index can be concealed by cancellation.
The actual exterior-power identity gives
\[
 \begin{aligned}
 \exp\!\left(\sum_{k\ge1}\frac{L_k}{k}t^k\right)
 &=\prod_{q=0}^4
       \det\bigl(I-t\Lambda^qT\bigr)^{(-1)^{q+1}}\\
 &=\prod_{n\ge1}(1-t^n)^{-a_n}.
 \end{aligned}                                  \tag{C.MS.7}\label{cdp:eq:C.MS.7}
\]
Every eigenvalue of every exterior power has
order dividing $M$. Hence the first product
factors into cyclotomic polynomials $\Phi_d(t)$
with $d\mid M$. The identity
$\Phi_d(t)=\prod_{n\mid d}(1-t^n)^{\mu(d/n)}$
for $d>1$, with the harmless sign at $d=1$,
expresses it using only indices dividing $M$.
Finite products of the factors $1-t^n$ have
unique integer exponent vectors: the coefficient
at the smallest differing degree in their
formal logarithms detects any difference.
Thus $a_n=0$ for every $n\nmid M$, including
all possible cancellations. If $T$ has an
eigenvalue one, every $L_k$ is zero and so is
the signature. Conversely, without that
eigenvalue, the real eigenvalues are $-1$
and the others occur in complex-conjugate
root-of-unity pairs. Their factors in
$\det(I-T)$ are respectively $2$ and
$|1-\lambda|^2$, all positive. Thus $a_1=L_1>0$.
The same factorization proves $L_k\ge0$ for
every $k$.

For \eqref{cdp:eq:C.MS.4}, choose the resolution with centres
over the singular locus of the original reduced
support. Its first exceptional multiplicities
are at least $2m_0$: at such a centre either
one support component has order at least two
or at least two components occur. All later
exceptional multiplicities are sums of the
orders of the components containing their
centres. If the generic point of a centre is
contained in an existing exceptional component,
that component already contributes at least
$2m_0$. Otherwise the map is locally an
isomorphism at its generic point and the
original reduced support still has order at
least two there. Thus this case has the same
lower bound, including successive resolutions
over disjoint original singular points.
The open stratum of a strict
transform is exactly the specified
$S_i^\circ$. In \eqref{cdp:eq:C.MS.6} for $k=m_0$, only
strict transforms of multiplicity $m_0$
contribute; this proves \eqref{cdp:eq:C.MS.4}.

If $m_0>1$, the resolved divisor has no
multiplicity-one component. Hence $a_1=0$,
so $T$ has eigenvalue one and the entire
signature vanishes. For an irreducible
fibre $mS$, \eqref{cdp:eq:C.MS.4} is
$m e(S_{\rm reg})=L_m$. When $m>1$ it
is zero by the preceding argument. When
$m=1$, its possible values are the
$a_1$ entries of the table, exactly the
set in \eqref{cdp:eq:C.MS.5}.

We now realize every row by an actual compact
family. The fifteen products of two degree-two
blocks are obtained from products of elliptic
curves. The five blocks
\[
                 \Phi_1^2,\quad\Phi_2^2,\quad
                 \Phi_3,\quad\Phi_4,\quad\Phi_6
\]
come respectively from the identity, negation,
and multiplication by a primitive third, fourth
or sixth root of unity on an elliptic curve
with the corresponding triangular or square
lattice. Taking the product automorphism gives
exactly the indicated product on integral
first cohomology.

For $\Phi_5$, let $A$ be the Jacobian of
the smooth genus-two curve
\[
                           C:\ y^2=x^5-1.
\]
Its double cover of $\PP^1$ has six simple
branch points, including infinity, so the
genus assertion follows from Riemann--Hurwitz.
The automorphism $(x,y)\mapsto(\zeta_5x,y)$
fixes infinity and acts on the basis
$dx/y,x\,dx/y$ of holomorphic differentials
with eigenvalues $\zeta_5,\zeta_5^2$.
The Abel--Jacobi map based at infinity
intertwines this automorphism with its induced
automorphism of $A$ and induces the actual
isomorphism on integral $H^1$. Including the
conjugate eigenvalues gives $\Phi_5$.
Composing the Jacobian automorphism with
negation gives order ten and characteristic
polynomial $\Phi_{10}$. Here negation is also
the map induced by the curve's hyperelliptic
involution: the divisor of $x-x(p)$ gives
$[p-\infty]+[\iota(p)-\infty]=0$ on the
Jacobian, verifying that comparison explicitly.

For the remaining two rows take
$A=E_i\times E_i$, where
$E_i=\C/(\Z+i\Z)$. The Gaussian-integer
matrices
\[
 G_8=\begin{pmatrix}0&i\\1&0\end{pmatrix},
 \qquad
 G_{12}=\begin{pmatrix}0&1\\1&i\end{pmatrix}
 \tag{C.MS.8}\label{cdp:eq:C.MS.8}
\]
define actual automorphisms of this torus.
One has $G_8^2=iI$, and its holomorphic
eigenvalues are $\zeta_8,\zeta_8^5$.
The characteristic polynomial of $G_{12}$
over $\C$ is $x^2-ix-1$, with roots
$(i+\sqrt3)/2=\zeta_{12}$ and
$(i-\sqrt3)/2=\zeta_{12}^5$.
Together with the conjugate eigenvalues,
these give the integral characteristic
polynomials $\Phi_8$ and $\Phi_{12}$ and
the exact orders eight and twelve.
Both matrices preserve the displayed lattice;
their determinants are Gaussian units.
This accounts for all nineteen rows with
finite-order automorphisms $g$ of projective
two-tori $A$.

Let $M$ be the order of the selected $g$.
Take the projective finite quotient
\[
 X=(A\times\PP^1)/
       \langle(a,z)\mapsto(g(a),\zeta_Mz)\rangle,
 \qquad f([a,z])=z^M.                           \tag{C.MS.9}\label{cdp:eq:C.MS.9}
\]
The action is free over $z\ne0,\infty$.
Over a nonzero finite base value the $M$
slices are permuted transitively, and their
quotient is the original connected torus $A$.
The actual positive meridian lifts from $z$
to $\zeta_Mz$, so under this identification
the monodromy is $g^{-1}$. Inverting all
the eigenvalues preserves each cyclotomic
polynomial and hence the displayed signature.
Resolve the projective quotient and then its
two central supports by projective blowups
over $0$ and $\infty$, using
\cite[Theorem~1.6]{BierstoneMilman1997}.
The resulting threefold is smooth and projective,
with SNC supports, and the punctured family is
unchanged. Its Stein factor has generic
degree one over the normal curve $\PP^1$,
so every fibre is connected. When $M=1$
the construction is simply the product,
which already has zero-Euler fibres.

For the remaining cases, each of the two
possibly singular supports has even Euler
number by \ref{cdp:st:C.2.UN.1}. Apply the exact vertical
blowups of \ref{cdp:st:C.2.UN.4} to make each number
zero. A point blowup changes it by $+2$;
a point blowup followed by the blowup of
a smooth plane quartic in its exceptional
$\PP^2$ changes it by $-2$. These
operations are projective and are performed
only over the selected value, so they
preserve the actual punctured family and
the signature \eqref{cdp:eq:C.MS.2}.

They preserve SNC supports as well. Choose
the point on a single component away from
the others. Its exceptional plane meets
the old strict transform along a line,
and choose the quartic transverse to that
line. Near any of their four intersection
points write the map as $t=(ru)^m$,
with exceptional component $r=0$, old
component $u=0$, and quartic centre
$(r,v)=(0,0)$. Its blowup charts
$v=rw$ and $r=vw$ give respectively
$t=(ru)^m$ and $t=(vwu)^m$.
Thus every resulting component is still
a coordinate hyperplane; away from these
four points the same assertion uses only
two coordinates. These explicit charts
verify the claimed SNC zero-Euler
realization of every row.
After the Euler number is zero, one additional
$+2$ operation and one $-2$ operation over the
same value leave it zero and introduce
exceptional components. Performing this at
both $0$ and $\infty$ ensures precisely two
reducible singular fibres, including the
identity-monodromy row. All other fibres
remain the original smooth tori.

The exact computation is saved in
\path{code/D_monodromy_signatures.py}, its
output, and
\path{records/D_monodromy_signatures.csv}.
It verifies every determinant and exterior-power
zeta identity over the integers or $\Q(t)$,
and the M\"obius inversions use exact integers.
The matrix-to-signature map itself forgets
information: for example $I_4$ and
$I_4+E_{12}$ have the same characteristic
polynomial and zero signature, while their
minimal polynomials differ. The table therefore
records precisely the stated invariant and
does not supply an analytic isomorphism
classification of fillings.
\end{proof}


